% Ampoule à décanter après extraction du limonène par le cyclohexane : position des phases (corrigé du devoir n° 2, exercice 3)
\begin{tikzpicture}[x=1cm,y=1cm]
\def\contour{(-0.16,4.1) -- (-0.16,3.75) .. controls (-0.7,3.45) and (-1.05,2.6) .. (-0.75,1.6)
  .. controls (-0.55,1.05) and (-0.12,0.75) .. (-0.07,0.55) -- (-0.07,-0.9) -- (0.07,-0.9) -- (0.07,0.55)
  .. controls (0.12,0.75) and (0.55,1.05) .. (0.75,1.6) .. controls (1.05,2.6) and (0.7,3.45) .. (0.16,3.75) -- (0.16,4.1)}
\begin{scope}
  \clip \contour -- cycle;
  \fill[pcEau!70!cyan!30] (-2,0.4) rectangle (2,1.95);
  \fill[yellow!35] (-2,1.95) rectangle (2,3.05);
\end{scope}
\draw[pcSurface] (-0.93,1.95) -- (0.93,1.95);
\draw[pcVerre] \contour;
\filldraw[pcAppareil] (-0.2,4.1) rectangle (0.2,4.38);
\filldraw[pcAppareil] (-0.17,-0.25) rectangle (0.17,0.05);
\filldraw[pcAppareil] (0.17,-0.15) rectangle (0.45,-0.05);
\draw[otInk!60,line width=1.1pt] (-1.15,1.25) -- (1.15,1.25);
\draw[otInk!60,line width=1.1pt] (1.15,1.25) -- (2.0,1.25) (2.0,-2.2) -- (2.0,4.6);
\pic at (0,-2.2) {becher={largeur=1.3,hauteur=1.1,niveau=0}};
\pcNote[west]{(0.55,2.6)}{(2.4,3.1)}{phase organique (dessus) :\\cyclohexane et limonène ; $d=0{,}78$}
\pcNote[west]{(0.55,1.4)}{(2.4,1.7)}{phase aqueuse (dessous) :\\eau ; $d=1{,}00$}
\pcNote[east]{(-0.2,4.25)}{(-1.2,4.4)}{bouchon}
\pcNote[east]{(-0.17,-0.1)}{(-1.2,-0.4)}{robinet}
\end{tikzpicture}
\par\smallskip\legende{Le cyclohexane, moins dense que l'eau, forme la phase supérieure : on élimine d'abord la phase aqueuse par le robinet.}
